%!TEX program = xelatex
\documentclass[cn,blue,11pt,normal]{elegantnote}

\title{高中物理题集}

\author{时胜举}
\institute{图灵教育}

\version{2.50}
\date{\jintian}

%%=============================================
%%================= 文档设置 ===================

%%================= 宏包调用 ===================
\usepackage{tikz}
\usepackage[version=4]{mhchem} % 化学公式
\usepackage{array,tabularx}
\usepackage{tabularray}
\usepackage{makecell,multirow,diagbox}
\usepackage{pifont}
\usepackage{xfrac} % 分数格式：斜分数\sfrac
\usepackage{multicol} % 分栏

%-------------- 自定义宏包 ---------------------
\usepackage{sapientiabook-choices}
\usepackage{sapientiabook-question}
%\usepackage{sapientiabook-answer}
\usepackage[show-answer]{sapientiabook-answer} % show-answer 显示答案

%%=============== 新建命令 =====================
\newcommand{\lkf}[1]{\noindent\text{\kaishu\color{cyan}（#1）}}
\newcommand{\jintian}{\number\year 年 \number\month 月 \number\day 日} % 中文年月日
\newcommand{\rmnum}[1]{\romannumeral #1} % 小写罗马数字
\newcommand{\Rmnum}[1]{\uppercase\expandafter{\romannumeral #1\relax}} % 大写罗马数字
\newcommand{\lunl}[1]{\underline{\hbox to #1{}}}%%  可输入长度的下划线
\newcommand{\lminiindent}[1]{\makebox[-1em]{#1\qquad}} % 缩进

\NewDocumentCommand\nceeID{o+mm}
	{%
		\IfNoValueTF{#3}%
		{\noindent\text{\kaishu\color{cyan}（#2，#3）}}%
		{\noindent\text{\kaishu\color{cyan}（#2，#3，#1分）}}%
	}
%%================ 全局设置 ====================
%-------------- 双线页眉的设置 ------------------
\makeatletter %双线页眉
\def\headrule{{\if@fancyplain\let\headrulewidth\plainheadrulewidth\fi %
\color{sakura}\hrule\@height 1.0pt \@width\headwidth\vskip1pt %上面线为1pt粗
\color{sakura}\hrule\@height 0.5pt\@width\headwidth%下面0.5pt粗
\vskip-2\headrulewidth\vskip-1pt} %两条线的距离1pt
\vspace{6mm}} %双线与下面正文之间的垂直间距
\makeatother
%----------------------------------------------
\pagestyle{fancy-note}

%%=============================================

\begin{document}
\markboth{高中物理\(\cdot\)题集}{festina lente}

%%-------------- 力学 -----------------------
	\input{mechanics/section_01}	%% 物理史、北京
	\input{mechanics/section_02}	%% 运动学
	\input{mechanics/section_03}	%% 相互作用
	\input{mechanics/section_04}	%% 动力学
	\input{mechanics/section_05}	%% 传送带与板块
	\input{mechanics/section_06}	%% 曲线
	\input{mechanics/section_07}	%% 引力
	\input{mechanics/section_08}	%% 功能
	\input{mechanics/section_09}	%% 动量
	\input{mechanics/section_10}	%% 振动与波

%%------------------- 电磁学 ---------------
	\input{electromagnetism/section_01}		%%电场
	\input{electromagnetism/section_02}		%%磁场
	\input{electromagnetism/section_03}		%%电磁场复合
	\input{electromagnetism/section_04}		%%电磁感应

%%---------------- 热学&光学&近物 -------------
	\input{thermology_optics_modern_physics/section_01}		%%热学
	\input{thermology_optics_modern_physics/section_02}		%%光学
	\input{thermology_optics_modern_physics/section_03}		%%原子物理
	
%%-------------- 实验  ----------------------
	\input{experiment/section_01}
	\input{experiment/section_02}
	\input{experiment/section_03}
\end{document}